\subsubsection{Subproceso 1: Acreditación y Pre-requisitos}

\begin{itemize}
    \item \textbf{Estudiante Regular:} Realiza el pago del autoseguro estudiantil, devuelve o paga cualquier material pendiente en su Facultad y revisa en los portales Intralu e Intranet DIRCE que sus requisitos figuren al día junto con su turno de matrícula asignado.

    \item \textbf{Dirección de Registro Central y Estadística (DIRCE / ORCE):} Consolida las notas del ciclo anterior, procesa las convalidaciones enviadas por las Facultades, calcula el promedio ponderado de los dos últimos semestres, asigna los turnos por orden de mérito y habilita la ficha de matrícula en el sistema SIGA-ORCE.

    \item \textbf{Comisión de Matrícula y Dirección de Escuela Profesional:} Coordinan con la DIRCE el cronograma de turnos de la Facultad, tramitan los cambios de plan de estudios y convalidaciones previas, y envían la lista de alumnos próximos a egresar para el turno de matrícula preferencial.

    \item \textbf{Unidad de Tesorería / Caja UNI y Dependencias de Facultad (Bibliotecas y Laboratorios):} Registran el pago del autoseguro estudiantil y reportan si el alumno mantiene deudas de libros, equipos o materiales de enseñanza en su Facultad.

    \item \textbf{Jefatura de Tutoría y Docentes Tutores:} Brindan la consejería previa obligatoria a los estudiantes con cursos desaprobados dos o tres veces y autorizan la activación de su código antes de que empiece la matrícula regular.
\end{itemize}
